% Planning van één labo. Gebruikt door de syllabus van dit labo en door de syllabus voor labo 3-5.
\textbf{Doel:} data aggregeren met \code{groupby}, tabellen koppelen met
\code{merge}, herschikken met \code{melt} en \code{pivot}, en eenvoudige
kenmerken van een tijdreeks berekenen.

\begin{schedule}
13:00--13:15 & Herhaling: een kapotte cel herstellen & Een kleine opkuis-pipeline met twee fouten repareren. \\
13:15--14:10 & \code{groupby} + \code{agg} & Statistieken per groep, leeftijdscategorieën met \code{pd.cut}, meerdere benoemde aggregaties tegelijk. \\
14:10--14:55 & \code{merge} & Inner en left joins; waarom een dubbele atleet twee keer meetelt, en waarom je dubbels verwijdert vóór het koppelen. \\
14:55--15:10 & Pauze & \\
15:10--15:50 & \code{melt} en \code{pivot} & Lange en brede vorm; vooruitgang tussen twee testmomenten. \\
15:50--16:50 & Tijdreeksen & Hartslagzones als \% van $HR_{max}$ (220 $-$ leeftijd), tijd per zone, \code{rolling}, \code{resample}. \\
16:50--17:30 & Minicase per twee & Een trainingsrapport: één rij per atleet. \\
\end{schedule}

\textbf{Oefeningen:} (0) herhaling: de kapotte cel herstellen;
(1) groeperen per geslacht en leeftijdscategorie; (2) wat heeft het
verwijderen van de dubbels gekost?; (3) vooruitgang in spronghoogte;
(4) minuten in zone 4--5 per atleet; \textbf{minicase:}
trainingsrapport met de minuten in zone 4--5, de maximale snelheid en de
afgelegde afstand per atleet.
